\documentclass[12pt,a4paper]{article}

\usepackage[utf8]{inputenc}
\usepackage[T1]{fontenc}
\usepackage{lmodern}
\usepackage[margin=2.5cm]{geometry}
\usepackage{graphicx}
\usepackage{hyperref}
\usepackage{listings}
\usepackage{xcolor}
\usepackage{booktabs}
\usepackage{enumitem}
\usepackage{fancyhdr}
\usepackage{titlesec}

\hypersetup{
    colorlinks=true,
    linkcolor=blue!70!black,
    urlcolor=blue!70!black
}

\lstset{
    basicstyle=\ttfamily\small,
    backgroundcolor=\color{gray!8},
    frame=single,
    framerule=0.4pt,
    rulecolor=\color{gray!40},
    breaklines=true,
    breakatwhitespace=true,
    tabsize=4,
    showstringspaces=false,
    numbers=none,
    xleftmargin=1em,
    framexleftmargin=0.5em
}

\pagestyle{fancy}
\fancyhf{}
\fancyhead[L]{\small CSE344 -- System Programming}
\fancyhead[R]{\small Homework \#2}
\fancyfoot[C]{\thepage}

% Boş etiket {} kullanmayın: bölüm numarası yazılmaz ve içindekiler boş kalabilir.
\titleformat{\section}{\large\bfseries}{\thesection\quad}{0em}{}[\titlerule]
\titleformat{\subsection}{\normalsize\bfseries}{\thesubsection\quad}{0em}{}
\setcounter{tocdepth}{3}

% Görsel klasörü (çoğul: screenshots). Tek yerden değiştirmek için:
\newcommand{\ScreenshotDir}{screenshots}

% ------------------------------------------------------------------
% Mandatory test screenshots: place files under screenshots/ (\ScreenshotDir)
% Run: pdflatex report.tex   (from the same directory as report.tex)
% Tries basenames #1, #2, #3 with .png then .pdf (see screenshots/README.txt)
% Note: avoid # inside \texttt in macro bodies (breaks \@argdef); use plain text for hints.
% ------------------------------------------------------------------
% Her adım için sırayla dener: #1 ve #2 ve #3 taban adları; .png .jpg .jpeg .pdf
% Linux'ta dosya adı büyük/küçük harfe duyarlı olabilir (ör. .PNG).
\newcommand{\MandatoryFigThree}[4]{%
\begin{figure}[htbp]
\centering
\IfFileExists{\ScreenshotDir/#1.png}{\includegraphics[width=0.95\textwidth]{\ScreenshotDir/#1.png}}{%
\IfFileExists{\ScreenshotDir/#1.jpg}{\includegraphics[width=0.95\textwidth]{\ScreenshotDir/#1.jpg}}{%
\IfFileExists{\ScreenshotDir/#1.jpeg}{\includegraphics[width=0.95\textwidth]{\ScreenshotDir/#1.jpeg}}{%
\IfFileExists{\ScreenshotDir/#1.pdf}{\includegraphics[width=0.95\textwidth]{\ScreenshotDir/#1.pdf}}{%
\IfFileExists{\ScreenshotDir/#1.PNG}{\includegraphics[width=0.95\textwidth]{\ScreenshotDir/#1.PNG}}{%
\IfFileExists{\ScreenshotDir/#2.png}{\includegraphics[width=0.95\textwidth]{\ScreenshotDir/#2.png}}{%
\IfFileExists{\ScreenshotDir/#2.jpg}{\includegraphics[width=0.95\textwidth]{\ScreenshotDir/#2.jpg}}{%
\IfFileExists{\ScreenshotDir/#2.jpeg}{\includegraphics[width=0.95\textwidth]{\ScreenshotDir/#2.jpeg}}{%
\IfFileExists{\ScreenshotDir/#2.pdf}{\includegraphics[width=0.95\textwidth]{\ScreenshotDir/#2.pdf}}{%
\IfFileExists{\ScreenshotDir/#2.PNG}{\includegraphics[width=0.95\textwidth]{\ScreenshotDir/#2.PNG}}{%
\IfFileExists{\ScreenshotDir/#3.png}{\includegraphics[width=0.95\textwidth]{\ScreenshotDir/#3.png}}{%
\IfFileExists{\ScreenshotDir/#3.jpg}{\includegraphics[width=0.95\textwidth]{\ScreenshotDir/#3.jpg}}{%
\IfFileExists{\ScreenshotDir/#3.jpeg}{\includegraphics[width=0.95\textwidth]{\ScreenshotDir/#3.jpeg}}{%
\IfFileExists{\ScreenshotDir/#3.pdf}{\includegraphics[width=0.95\textwidth]{\ScreenshotDir/#3.pdf}}{%
\IfFileExists{\ScreenshotDir/#3.PNG}{\includegraphics[width=0.95\textwidth]{\ScreenshotDir/#3.PNG}}{%
\fbox{\parbox{0.9\textwidth}{\centering\vspace{1em}\small
\textit{Missing image.} Put the screenshot in \texttt{\ScreenshotDir/} next to \texttt{report.tex}.
Try basenames \texttt{\detokenize{#1}}, \texttt{\detokenize{#2}}, or \texttt{\detokenize{#3}}
with extension \texttt{.png}, \texttt{.jpg}, \texttt{.jpeg}, or \texttt{.pdf}.
Re-run \texttt{pdflatex} twice after adding files.\vspace{1em}}%
}}}}}}}}}}}}}}}%
}%
\caption{#4}%
\end{figure}%
}

\begin{document}

% ============================================================
% COVER PAGE
% ============================================================
\begin{titlepage}
    \centering
    \vspace*{3cm}

    {\Large\textbf{Gebze Technical University}\\
    Department of Computer Engineering}

    \vspace{2cm}

    {\LARGE\textbf{CSE344 -- System Programming\\Homework \#2}}

    \vspace{1cm}

    {\Large Signal Handling \& Process Management with \texttt{fork()}}

    \vspace{3cm}

    \begin{tabular}{rl}
        \textbf{Name:}       & [Salih Cengiz] \\
        \textbf{Student ID:} & [220104004007] \\
        \textbf{Course:}     & CSE344 -- System Programming \\
        \textbf{Date:}       & March 2026 \\
    \end{tabular}

    \vfill
\end{titlepage}

% İçindekiler: iki kez pdflatex gerekir (birinci .toc yazar, ikinci doldurur).
\tableofcontents
\newpage

% ============================================================
% 1. INTRODUCTION
% ============================================================
\section{Introduction}

This report describes the design and implementation of \texttt{procSearch}, a multi-process file search tool written in C for Linux. The program accepts a root directory and a filename pattern, forks a configurable number of worker processes, and assigns each worker a partition of the directory tree to search recursively. Workers report results to the parent through signals, and the parent produces a final tree-formatted output with a summary.

The program is built with a Makefile, compiles with \texttt{-Wall} with zero warnings, and runs on Debian 11 (64-bit). All inter-process communication relies on POSIX signals (\texttt{SIGUSR1}, \texttt{SIGINT}, \texttt{SIGTERM}, \texttt{SIGCHLD}) and temporary files for match data transfer.

\subsection{Key Features}

\begin{itemize}[noitemsep]
    \item Custom pattern matching with \texttt{+} operator (one or more of the preceding character), implemented from scratch without any regex library
    \item Round-robin directory partitioning among worker processes
    \item Full signal handling with \texttt{sigaction()} and async-signal-safe handlers
    \item Tree-formatted output with per-worker match summary
    \item Zombie process prevention via \texttt{SIGCHLD} handler
    \item Graceful shutdown on \texttt{SIGINT} (Ctrl+C)
\end{itemize}

% ============================================================
% 2. DESIGN DECISIONS
% ============================================================
\section{Design Decisions}

\subsection{Overall Architecture}

The program follows a parent-worker model:

\begin{lstlisting}
Parent Process
  |
  |-- fork() --> Worker 0  (assigned: alpha, delta, ...)
  |-- fork() --> Worker 1  (assigned: beta, ...)
  |-- fork() --> Worker 2  (assigned: gamma, ...)
  |
  |-- wait for SIGUSR1 from all workers
  |-- collect results from temp files
  |-- print tree output and summary
\end{lstlisting}

The parent is responsible for argument parsing (\texttt{getopt}), directory partitioning, forking workers, waiting for completion, and producing the final output. Workers perform the actual recursive directory traversal and pattern matching.

\subsection{Directory Partitioning}

The \texttt{read\_subdirectories()} function reads the immediate children of the root directory using \texttt{opendir()}/\texttt{readdir()}. Only directories (checked via \texttt{stat()} and \texttt{S\_ISDIR}) are collected. These subdirectories are distributed among workers using round-robin assignment:

\begin{lstlisting}
subdirectory[i] --> worker[i % num_workers]
\end{lstlisting}

Two edge cases are handled:

\begin{itemize}[noitemsep]
    \item \textbf{Fewer subdirectories than workers:} The worker count is reduced to match the subdirectory count, and a notice is printed.
    \item \textbf{No subdirectories at all:} The parent searches the root directory directly via \texttt{parent\_direct\_search()} without forking any workers.
\end{itemize}

\subsection{Pattern Matching Algorithm}

The \texttt{match\_pattern()} / \texttt{match\_pattern\_at()} functions implement a custom pattern matching engine supporting the \texttt{+} operator. The algorithm performs substring matching (the pattern can match anywhere within the filename) and is case-insensitive (using \texttt{tolower()}).

The \texttt{+} operator applies to the character immediately to its left and means ``one or more occurrences.'' For example, \texttt{rep+ort} matches \texttt{report}, \texttt{repport}, \texttt{reppport}, etc.

A key design decision was implementing \textbf{backtracking} for the \texttt{+} operator. When a \texttt{c+} group is encountered:

\begin{enumerate}[noitemsep]
    \item All consecutive occurrences of \texttt{c} in the filename are consumed (greedy).
    \item The remaining pattern is tested against the remaining filename.
    \item If the test fails, the number of consumed characters is reduced by one and retried.
    \item This continues until either a match is found or the minimum of 1 repetition is reached.
\end{enumerate}

This backtracking is essential for patterns like \texttt{er+ro+r} matching \texttt{error}, where \texttt{r+} must not consume both \texttt{r}s so that the subsequent literal \texttt{r} in the pattern can also match.

The backtracking is implemented via recursion in \texttt{match\_pattern\_at()} --- when a \texttt{+} group is encountered, the function calls itself with the remaining pattern and each possible consumption length.

\subsection{Signal Handling Strategy}

All signal handlers are installed using \texttt{sigaction()} (not \texttt{signal()}) for reliable, portable behavior. The handlers are minimal and async-signal-safe:

\begin{table}[h]
\centering
\begin{tabular}{@{}llp{7.5cm}@{}}
\toprule
\textbf{Signal} & \textbf{Handler} & \textbf{Behavior} \\
\midrule
\texttt{SIGUSR1} & \texttt{handle\_sigusr1} & Increments \texttt{finished\_worker\_count} (\texttt{volatile sig\_atomic\_t}) \\
\texttt{SIGINT}  & \texttt{handle\_sigint}  & Sets \texttt{sigint\_received} flag \\
\texttt{SIGTERM} & \texttt{handle\_sigterm} & Sets \texttt{sigterm\_received} flag (worker-side) \\
\texttt{SIGCHLD} & \texttt{handle\_sigchld} & Calls \texttt{waitpid(-1, \&status, WNOHANG)} in a loop, records PID/status \\
\bottomrule
\end{tabular}
\end{table}

\subsubsection{Async-Signal-Safety Measures}

\begin{itemize}[noitemsep]
    \item \textbf{No \texttt{printf()}, \texttt{malloc()}, or \texttt{free()} inside any handler.} All handlers only modify \texttt{volatile sig\_atomic\_t} variables or call \texttt{waitpid()} (which is async-signal-safe).
    \item \textbf{\texttt{errno} preservation:} The \texttt{SIGCHLD} handler saves and restores \texttt{errno} to avoid corrupting the main program's error state.
    \item \textbf{\texttt{SA\_RESTART} is intentionally not used} for \texttt{SIGUSR1} and \texttt{SIGCHLD} handlers. This ensures that \texttt{sleep()} in the parent's wait loop is interrupted when a signal arrives, allowing the loop condition to be re-evaluated immediately.
    \item \textbf{\texttt{SA\_NOCLDSTOP}} is set for \texttt{SIGCHLD} to avoid spurious signals when children are stopped (not terminated).
\end{itemize}

\subsubsection{Signal Flow}

\begin{lstlisting}
Worker finishes  --> SIGUSR1 --> Parent (increment counter)
User Ctrl+C      --> SIGINT  --> Parent (send SIGTERM to all workers)
Parent SIGTERM   --> SIGTERM --> Worker (stop, report partial results)
Worker exits     --> SIGCHLD --> Parent (reap zombie, record status)
\end{lstlisting}

\subsection{Worker-Parent Communication}

Workers communicate match results to the parent through temporary files:

\begin{itemize}[noitemsep]
    \item \textbf{Match data:} \texttt{/tmp/procSearch\_<PID>.tmp} --- one line per match in \texttt{path|size|worker\_pid} format
    \item \textbf{Scan count:} \texttt{/tmp/procSearch\_<PID>\_info.tmp} --- contains the number of files scanned
\end{itemize}

This approach avoids the complexity of pipes or shared memory while being safe across process boundaries. The parent reads these files after all workers have finished and deletes them afterward.

\subsection{Wait Loop Design}

The parent waits for all workers using a \texttt{while} loop with \texttt{sleep(1)}:

\begin{lstlisting}[language=C]
while (finished_worker_count < total_worker_count &&
       reaped_count < total_worker_count &&
       !sigint_received) {
    sleep(1);
}
\end{lstlisting}

The loop checks two independent completion indicators:

\begin{itemize}[noitemsep]
    \item \texttt{finished\_worker\_count}: incremented by \texttt{SIGUSR1} handler
    \item \texttt{reaped\_count}: incremented by \texttt{SIGCHLD} handler
\end{itemize}

This dual-condition approach provides robustness: even if a \texttt{SIGUSR1} signal is lost (e.g., due to signal coalescing), the \texttt{SIGCHLD}-based count ensures the parent eventually exits the loop.

% ============================================================
% 3. CHALLENGES AND SOLUTIONS
% ============================================================
\section{Challenges and Solutions}

\subsection{Challenge 1: Parent Hanging in Wait Loop}

\textbf{Problem:} The parent process would hang indefinitely in the wait loop, even after all workers had finished.

\textbf{Root Cause:} \texttt{SA\_RESTART} was set for \texttt{SIGUSR1} and \texttt{SIGCHLD} handlers. When a signal arrived during \texttt{sleep(1)}, the kernel automatically restarted the \texttt{sleep()} call instead of returning, preventing the loop from re-evaluating its condition.

\textbf{Solution:} Removed \texttt{SA\_RESTART} from \texttt{SIGUSR1} and \texttt{SIGCHLD} signal configurations. Now \texttt{sleep()} is interrupted by any signal, the function returns early, and the loop condition is checked immediately.

\subsection{Challenge 2: Greedy Pattern Matching Without Backtracking}

\textbf{Problem:} The \texttt{+} operator consumed all matching characters greedily with no way to try fewer matches. For example, \texttt{er+ro+r} failed to match \texttt{error} because \texttt{r+} consumed both \texttt{r} characters, leaving none for the subsequent literal \texttt{r} in the pattern.

\textbf{Solution:} Implemented backtracking in \texttt{match\_pattern\_at()}. When a \texttt{c+} group is encountered, the function first finds the maximum number of matches, then tries each possible consumption count from maximum down to 1 by recursively calling itself with the remaining pattern.

\subsection{Challenge 3: Duplicate Notice Messages After fork()}

\textbf{Problem:} The notice message was printed multiple times --- once by the parent and once by each forked child.

\textbf{Root Cause:} \texttt{stdout} is fully buffered when redirected to a pipe. The \texttt{printf()} output remained in the buffer when \texttt{fork()} was called, and each child inherited a copy of the buffer. When each child exited, its copy was flushed, producing duplicate output.

\textbf{Solution:} Added \texttt{fflush(stdout)} immediately after printing the notice message and before any \texttt{fork()} call.

\subsection{Challenge 4: Stack Overflow Risk}

\textbf{Problem:} Several large arrays were allocated on the stack, risking stack overflow.

\textbf{Solution:} Large arrays in \texttt{main()} and helper functions were declared \texttt{static}, moving them from the stack to the data segment. Additionally, \texttt{MAX\_DIRS\_PER\_WORKER} was set to 512 to limit per-worker memory footprint.

\subsection{Challenge 5: Collecting Exit Status for Zero-Match Workers}

\textbf{Problem:} Workers with 0 matches called \texttt{exit(0)}. The fallback \texttt{waitpid()} loop used \texttt{match\_count > 0} to determine if a worker was already collected, causing zero-match workers to be redundantly processed.

\textbf{Solution:} Added a \texttt{collected} flag to the \texttt{WorkerInfo} struct. This flag is set when a worker's exit status is processed, regardless of the match count value.

% ============================================================
% 4. MANDATORY TEST SCENARIO (screenshots)
% ============================================================
\section{Mandatory Test Scenario}

This section follows the assignment PDF \textit{Mandatory Test Scenario}. Screenshots are loaded from the \texttt{\ScreenshotDir/} folder next to \texttt{report.tex}. Compile with:

\begin{verbatim}
cd <HW2_directory>
pdflatex report.tex
pdflatex report.tex
\end{verbatim}

For each step, the first existing file among \texttt{StepN}, \texttt{stepN}, or \texttt{step-0N} with extension \texttt{.png} or \texttt{.pdf} is included. Rename your captures to match, or adjust the macro arguments in the source.

\subsection{Step 1: Test directory tree under \texttt{/tmp/crawltest}}

Create the directory structure and files as specified (e.g.\ \texttt{mkdir}, \texttt{echo}, \texttt{touch}). Capture terminal output showing the tree or \texttt{find}/\texttt{ls} listing.

\MandatoryFigThree{Step1}{step1}{step-01}{Step 1: Directory tree and files under \texttt{/tmp/crawltest} as required by the assignment.}

\subsection{Step 2: Basic search with \texttt{rep+ort}}

Command (adjust pattern per instructor revision if needed):

\begin{verbatim}
./procSearch -d /tmp/crawltest -n 3 -f 'rep+ort'
\end{verbatim}

\MandatoryFigThree{Step2}{step2}{step-02}{Step 2: \texttt{procSearch} run showing live \texttt{MATCH} lines, tree output, and summary.}

\subsection{Step 3: Size filter \texttt{-s 15}}

\begin{verbatim}
./procSearch -d /tmp/crawltest -n 3 -f 'rep+ort' -s 15
\end{verbatim}

\MandatoryFigThree{Step3}{step3}{step-03}{Step 3: Only files with size $\geq$ 15 bytes appear in results.}

\subsection{Step 4: Pattern with no matches}

\begin{verbatim}
./procSearch -d /tmp/crawltest -n 3 -f 'xyz+123'
\end{verbatim}

\MandatoryFigThree{Step4}{step4}{step-04}{Step 4: Output shows \texttt{No matching files found.} and correct summary.}

\subsection{Step 5: Large directory, \texttt{SIGINT}, and zombie check}

\begin{verbatim}
./procSearch -d /usr/share/doc -n 4 -f 'rep+ort'
# While running: Ctrl+C, then:
ps -ef | grep procSearch | grep -v grep
\end{verbatim}

\MandatoryFigThree{Step5}{step5}{step-05}{Step 5: \texttt{SIGINT} shutdown message, clean termination, and/or \texttt{ps} showing no leftover/zombie \texttt{procSearch} processes.}

\subsection{Step 6: \texttt{stderr} / log evidence for signal-related behavior}

Example (redirect stderr while testing):

\begin{verbatim}
./procSearch -d /tmp/crawltest -n 3 -f 'rep+ort' 2> run.stderr
# Or strace (optional):
strace -f -e trace=signal -o strace.log ./procSearch ...
\end{verbatim}

\MandatoryFigThree{Step6-sdterr}{Step6}{step6}{Step 6: \texttt{stderr} or terminal capture relevant to \texttt{SIGINT} handling and/or unexpected worker termination (\texttt{SIGCHLD} path), as applicable to your run.}

% ============================================================
% 5. CONCLUSION
% ============================================================
\section{Conclusion}

The \texttt{procSearch} tool successfully implements all required features: custom pattern matching with backtracking, round-robin directory partitioning, multi-process search with \texttt{fork()}, comprehensive signal handling, and tree-formatted output. The program compiles with \texttt{-Wall} and zero warnings, correctly prevents zombie processes, and handles edge cases such as insufficient subdirectories and graceful shutdown via Ctrl+C. All signal handlers are async-signal-safe, using only \texttt{volatile sig\_atomic\_t} variables and async-signal-safe system calls.

\end{document}
